\chapter*{Annotation guidelines}
\label{cha:annotation_guide}

Loosely based on \cite{benikova-etal-2014-nosta}, \cite{ldc_2005}, \cite{muc7_ner}, \cite{smartdata_anno}


\section{Named Entity - Basic Definition}

Names are denominators for a single unique object in the real world. Only full noun phrases are considered as named entities. Pronouns and all other phrases can be ignored. 

Named entities can consist of more than one token. Phrases only partly containing names (e.g. `deutschlandweit), or adjectives referring to entities (`euklidisch') can be ignored.

Combined named entities (e.g. `Wien-Demo') should be tagged separately if possible (Wien[LOC]-Demo[ACTION]), or else tagged as the dominant entity type (here: [ACTION]) to avoid nested entities.

WHITESPACE: Not labelled

ONLY NOUNS: No verbs can be entities (Spaziergang vs spazieren gehen), dates and times can be represented by numbers and strings. 


%https://www-nlpir.nist.gov/related_projects/muc/proceedings/ne_task.html

\section{Definition of Different Entity Types}
Person (PER) - Named person. Person entities are limited to humans. A person must be a single
individual. Fictional people can be included, as long as they’re referred to by name (e.g. Harry Potter). 

Organization (ORG) - ORGANIZATION: named corporate, governmental, or other organizational entity. Organization entities are corporations, agencies, and other groups of people. Companies, agencies, institutions, newspapers, tv/radio channels, other media platforms; political parties; non-governmental organizations; political groups

Location (LOC) - Name of politically or geographically defined location (cities, provinces, countries, international regions, bodies of water, mountains, etc.). Location entities are geographical entities such as geographical areas and landmasses, bodies of water, and geological formations; furthermore they are geographical regions defined by political and/or social groups. Locations also include buildings, and addresses including streets, squares, and postal codes. Countries, cities, states; continents, districts, regions, streets and squares; shopping centers, parking lots, restaurants.

Date (DATE): Absolute or relative date expression. 

Time (TIME): Absolute time expressions, including times of day and time codes (e.g. for video urls). 

Action (ACTION): Specific type of action undertaken to express political beliefs, e.g. demonstrations, vigils, motorcades, walks in public. Specific action or event mentioned in a text, e.g. demonstrations, pickets, motorcade; (protest) walks; etc. - must not be unique entity, plural allowed

Frechheit - WNUT 15: %https://docs.google.com/document/d/12hI-2A3vATMWRdsKkzDPHu5oT74_tG0-PPQ7VN0IRaw/edit


% hier wird klar wie der shift von protest forms is: Achtung Demonstration wird Abgesagt ! Wir haben als Orga Team beschlossen , das wir keine Demos mehr anmelden werden ! Wir sagen klar Nein zu den Auflagen die uns vorgegeben werden . Deshalb , neuer Termin , wir fahren nach Plauen ! Lasst uns die Freie Sachsen und den Bürgermeisterkandidaten Thomas Kaden unterstützen ! 11062021 Raus auf die Straße


\begin{table}[]
\caption{Annotation scheme: Entity types}
\label{tab:entity_types}
\centering
\begin{adjustbox}{max width=\textwidth}
\begin{tabular}{p{0.25\linewidth}|p{0.65\linewidth}|p{0.65\linewidth}}

\textbf{Entity type} & \textbf{Description} & \textbf{Examples} \\
\hline 

PER &  & \\
\hline

LOC &  &  \\ 
\hline

ORG &  & \\ 
\hline

DATE & Absolute or relative dates & 23.09.2022 \\ 
\hline

TIME & Specific times of day, time codes & 13:10 \\ 

\hline

ACTION &  &  \\ 

\hline

\end{tabular}
\end{adjustbox}

\end{table}

\section{Specifications}

\subsection{Mixed-language messages}

Mixed-language where actual message is non-German language should be skipped. 

Examples:
We are angry and pissed off - Raus auf die Straße \@Demotermine!
Brilliant numbers in Carlisle ! Raus auf die Straßen @Demotermine

\subsection{Incomprehensible messages}
Incomprehensible message should be skipped. 

Example: 
Sterne-Restaurant Walliserkanne . Naudreck . Bersets faule Tricks . Uniarzt panikfrei . Herzmuskel .  Neuer Beitrag \#Walliserkanne Raus auf die Straßen \@Demotermine! Übersicht / Overview

\subsection{Special characters}
SPECIAL CHARS: Only label if in the middle of sequence representing an an entity

Examples:
Donnerstag[B-DATE] ,[I-DATE] 17.04.2020[I-DATE]

\subsubsection{Hashtags}
HASHTAGS - If hashtag -> include
Examples:
\#Berlin[B-ORG]
\#BW[B-ORG]
\#CH[B-ORG]
\#Hardy[B-PER] Groeneveld[I-PER]
\#Mutigmacher[B-ORG]n

\section{Specifications and Detailed Examples for Entity Types}

\subsection{Locations}
LOCATION
Alex, Nollendorfplatz
Norden
Hellersdorf
Berlin
Deutschland
Frankreich
Bundestag

No derived location-words: 
Kreuzberger[O] Nächte[O]
Combined locations
 e.g. streets or routes -> B-LOC, I-LOC, incuding special chars (e.g. , - . /)

Examples: 
Straße[B-LOC] des[I-LOC] 17.[I-LOC] Juni[I-LOC]
Bad[B-LOC] Oldesloe[I-LOC]
Elbe[B-LOC] Elster[I-LOC]
Bremen[B-LOC]-Vegesack[I-LOC]
Königstraße[B-LOC] Richtung[I-LOC] HBF[I-LOC]
Bernauer[B-LOC] Str[I-LOC] Richtung[I-LOC] Süden[I-LOC]
Marktoberdorf[B-LOC] -[I-LOC] Kaufbeuren[I-LOC]-Neugablonz[I-LOC]
Lutherstadt[B-LOC] Wittenberg[I-LOC]
D-67454[B-LOC] Haßloch[I-LOC]
Landeplatz[B-LOC] Grießkirchen[I-LOC]
Google[B-LOC] center[I-LOC]
Stuttgarter[B-LOC] Landtag[I-LOC]

ABER:
Berlin[B-LOC] Neptunbrunnen[B-LOC]


Combined location and date hashtag: 
\#b[B-LOC]2908[B-DATE]

Metaphorical locations:  not labeled
e.g. Berliner[O] Regenwald[O]

\subsection{Organizations}
ORGANIZATION
Bundesregierung[B-ORG
Gesundheitsamt[B-ORG
Französische[B-ORG] Polizei[I-ORG]
Facebook[B-ORG]
Telegram[B-ORG]
Google[B-ORG]
Bill[B-ORG]  Gates[I-ORG] Foundation[I-ORG]
ARD[B-ORG]
KSV[B-ORG] LiLi[I-ORG]
FC[B-ORG] St.[I-ORG] Pauli[I-ORG]
SGD[B-ORG]-Ultras[I-ORG]
Bündnis21[B-ORG] Hessen[I-ORG]
Polizei[B-ORG]
Polizeipräsidium[B-ORG] Pforzheim[I-ORG]


ORG vs. LOC? -> LOC 

Examples: 
Landtag[B-LOC]
Charité[B-LOC]
Charité[B-LOC] krankenhaus[I-LOC]



\subsection{Persons}
PERSON 
Scholz[B-PER]
Carolin[B-PER] Matthie[I-PER]
Anselm[B-PER] Lenz[I-PER]
Karl[B-PER]
Merkel[B-PER]s

\subsection{Time}
TIME 
Times of day: 
15:00[B-TIME]
17:00[B-TIME] Uhr[I-TIME]

Time codes after URLS, eg. www.example.com 00:30[B-TIME]

Relative time values: 
Mittag[B-DATE]
15:30[B-TIME] -[I-TIME] 19:30[I-TIME]

\subsection{Date}
DATE
2019[B-DATE]
13.12.21[B-DATE]
2.7.21[B-DATE]
31. [B-DATE] JULI[I-DATE]
01122020[B-DATE]
Pfingstmontag[B-DATE]
HEUTE[B-DATE]


Combined dates / periods
 -> B-DATE, I-DATE
Examples: 
17042021[B-DATE] –[I-DATE] 21042021[I-DATE]
28. [B-DATE]+ [I-DATE] 29.8. [I-DATE]


\subsection{Action}
ACTION
Angriff-B
Demo-B
Autokorso-B
Offenes-B Treffen-I
Aufzug
Zug
Plakate
Infostände
Informationstouren
Anti[B]-Corona[I]-Proteste[I]
Protestzug
Spaziergänge
Querdenker[B-ORG]-Demo[B-ACTION]
Korso[B-ACTION] Nord[I-ACTION]
\#kidsfreedomday[B-ORG]
Gedenken
Aktionswoche
Wahlkampf
Live-Stream

Not labeled, to imprecise:
Super Aktion[O] in Bönnigheim[B-LOC]


ACTION-LOCATION
Demo[B-ACTION]-[O]Berlin[B-LOC]
separately; unsure if they are NEs? 





hashtags encoding events via location and date, e.g. \#b20200809 or \#le1311 are not to be annotated, should be extracted via regex and treated separately; unsure if they are NEs? 

